\section{Applications and extensions}
\label{ch:applications}

Applications differ in their spatial margins, available modes, and empirical
flow construction, even though they use the same transport operators. The U.S.
freight exercise shows the complete economic-geography application. The urban
material that follows changes the spatial equilibrium while retaining the
recursive routing, modal, and congestion maps.

\subsection{U.S. freight}
\label{ch:rsue}

The freight application maps domestic and international value flows onto a
multimodal U.S. network and evaluates bidirectional road improvements. We
first define the exercise and its flow construction before reporting the
specification, results, decomposition, and provenance.

\subsubsection{Purpose}

The paper applies the statistic to the U.S. freight network
\citep{AllenFuchsWong2026}. The archived paper-replication result has
\GuideRSUENodeCount{} nodes: 228 domestic locations and six foreign regions.
The policy set contains \GuideRSUEDirectedArcCount{} directed Interstate road
arcs paired into \GuideRSUEPhysicalLinkCount{} bidirectional physical links.
Each experiment lowers the primitive road cost by one percent in both
directions of one physical link.

The exercise is a sufficient-statistic application. It does not estimate a
new 234-location structural model from raw microdata. The adapter combines
model-ready domestic trade and modal flow matrices, location shares, the road
network, and directional 2017 Census port trade. It then checks and constructs
the inputs required by the derivative.

\subsubsection{International flows}

The headline vintage preserves separate import and export directions from
public 2017 Census containerized vessel trade
\citep{CensusInternationalTradePort}. Port values are mapped to domestic nodes
and six foreign regions. The adapter first normalizes imports and exports
separately. It then averages the normalized U.S.-port margins across the two
directions and does the same for the foreign-region margins. Iterative
proportional fitting projects each directional matrix onto these common
margins while preserving its observed zero pattern. The routine requires a
maximum margin error below \(10^{-12}\); it fails if the observed support
cannot sustain the common margins or if the iteration does not converge.
Finally, the adapter assigns one-half of the normalized foreign-water modal
mass to each direction. This balancing step is part of the model adapter. It
does not assert that raw port-level imports and exports are equal.

The raw monthly API cache is not committed. The repository instead provides a
credential-safe downloader, port and foreign-region crosswalks, and the
aggregate derived overlay. It also records source metadata and hashes, along
with diagnostics for coverage, iteration counts, margin errors, total
variation, and node-level residuals.
The workflow reads the API key from the environment. Its tests and secret scan
verify that the tracked cache metadata, URLs, manifests, and outputs contain no
key.

\subsubsection{Headline specification}

The archived configuration uses:
\[
 \alpha=0.10,\quad \beta=-0.30,\quad \sigma=9,\quad
 \eta=1.099,\quad \gamma_{\mathrm{road}}=0.092.
\]
Subsection \ref{ch:data} documents the corresponding calibration sources:
\citet{AllenArkolakis2014,AllenArkolakis2020} and
\citet{FuchsWong2026Multimodal}.
It uses \code{ChoiceLogsum}, all observed transport modes in equilibrium,
edge-local road congestion, and no headline terminal congestion.
The configuration field retains the historical name \code{lambda\_road};
\(\gamma_{\mathrm{road}}\) is the theoretical congestion elasticity used in
this guide.
These parameters imply \(e=-0.5\) and \(\rho=-1.6\). The common factor
\(\rho\) scales all links equally. Traffic, pass-through, and the two
market-access multipliers generate the cross-link ranking. The sign of
\(\rho\) alone does not determine the welfare sign because the multiplier sum
is also signed; the archived baseline produces positive benefits.

The policy is a primitive road-cost improvement. The Traditional approach uses
the two-direction road traffic share. The full realized-cost effect changes
the cost users face after congestion. The Extended approach reports the full
primitive-cost effect and includes route, mode, congestion, and
spatial-equilibrium responses.

\subsubsection{Archived paper-replication results}

The current package reproduces the archived numerical reference:
\begin{table}[H]
\centering
\caption{Archived physical-link result summary}
\label{tab:rsue-summary}
\begin{tabular}{lr}
\toprule
Measure & Value \\
\midrule
Physical links & \GuideRSUEPhysicalLinkCount \\
Mean Traditional-approach elasticity & \GuideRSUEMeanTraditional \\
Mean full realized-cost effect & \GuideRSUEMeanRealized \\
Mean full primitive-cost effect & \GuideRSUEMeanExtended \\
Median primitive-cost effect & \GuideRSUEMedianExtended \\
Maximum primitive-cost effect & \GuideRSUEMaximumExtended \\
Pearson correlation with Traditional approach & \GuideRSUEPearson \\
Spearman rank correlation & \GuideRSUESpearman \\
\bottomrule
\end{tabular}
\end{table}

\implementationbox{These values pass the package's accounting,
algebraic, finite-difference, and reproducibility checks. Those checks verify
the baseline construction and the derivative calculation. They do not
validate the calibrated congestion, modal, or spatial responses against
untargeted empirical outcomes.}

The high correlations show that traffic tracks the Extended approach closely
across links.
They do not imply identical targeting. Only \GuideRSUETopTenOverlap{} links
overlap between the two top-ten lists in the archived result vintage. The
paper reports names only where endpoints can be assigned by point-in-polygon
matching to public 2018 Census CBSA geography.

\subsubsection{Signed decomposition}

To explain why the mean primitive-cost derivative differs from traffic, apply the
signed Hulten-gap identity in Subsection \ref{ch:decomposition} and average each
level contribution across physical links:
\[
\begin{array}{lr}
\text{Externality scale} & \GuideRSUEExternalityComponent\\
\text{Equilibrium propagation} & \GuideRSUEPropagationComponent\\
\text{Road congestion} & \GuideRSUERoadComponent\\
\text{Terminal congestion} & \GuideRSUETerminalComponent\\
\text{Primitive-cost pass-through} & \GuideRSUEPassThroughComponent .
\end{array}
\]
The externality and propagation terms have opposite signs, and each is large
in absolute value relative to their sum. They therefore nearly offset. The
full realized-cost mean exceeds the Traditional approach, whereas primitive-cost
pass-through lowers the headline primitive result. The decomposition reports
signed welfare levels because dividing these terms by the small net gap would
produce unstable shares.

\subsubsection{Interpretation}

The result ranks local model-implied benefits, not project net present values.
A high-ranked link may be expensive to build, face environmental constraints,
or conflict with distributional objectives. Conversely, a low-ranked capacity
improvement may be justified by safety or resilience benefits outside the
model. The appropriate use is to add the spatial-equilibrium benefit measure
to a broader appraisal, not to suppress the rest of the appraisal.

\subsubsection{Result provenance}

The guide asset script reads a public-safe aggregate replication record and
checks it against the declared configuration, input, source-result, and output
hashes in \path{replication/rsue/expected_summary.toml}. With restricted
results present, it also verifies the full directed and physical output files.
The archived vintage requires the stated network counts, the numerical gates,
a hash-bound finite-difference report, and deterministic outputs. Appendix
\ref{app:provenance} records the reproducibility boundary; the repository audit
contains the developer-level kernel lineage.

\subsubsection{Reproducing with external inputs}

Set the external data root and run the locked environment:
\begin{lstlisting}[language=bash]
export RSUE_DATA_ROOT=/path/to/hash-verified/rsue-inputs
julia --project=replication/rsue/environment \
  -e 'using Pkg; Pkg.instantiate()'
julia --project=replication/rsue/environment \
  replication/rsue/verify_choice_logsum_fd.jl \
  replication/rsue/rsue_paper_choice_edge_census_2017.toml \
  replication/rsue/verification/choice_logsum_rsue_fd.toml
julia --project=replication/rsue/environment \
  replication/rsue/build_paper_artifacts.jl
\end{lstlisting}
The manifest rejects files whose hashes differ from the declared data vintage.
There is no fallback to the legacy port matrix.

The freight application keeps the paper's economic-geography model. The
urban application instead changes the spatial residuals and welfare projection
while retaining the same route, modal, and congestion operators. The next
subsection describes its implementation and data requirements.
